\documentclass{article}
\usepackage{amsmath, amssymb, amsthm}
\usepackage{hyperref}
\usepackage{geometry}
\geometry{margin=1in}

\title{Erd\H{o}s Problem \#521}
\date{Last updated: 2025-08-31}
\author{Source: erdosproblems.com}

\begin{document}
\maketitle

\noindent\textbf{Status:} OPEN \quad \textbf{No prize}

\noindent\textbf{Tags:} analysis, polynomials, probability

\medskip

\noindent\textbf{Problem Statement:}

Let $(\epsilon_k)_{k\geq 0}$ be independently uniformly chosen at random from $\{-1,1\}$. If $R_n$ counts the number of real roots of $f_n(z)=\sum_{0\leq k\leq n}\epsilon_k z^k$ then is it true that, almost surely,\[\lim_{n\to \infty}\frac{R_n}{\log n}=\frac{2}{\pi}?\]




Erd\H{o}s and Offord \cite{EO56} showed that the number of real roots of a random degree $n$ polynomial with $\pm 1$ coefficients is $(\frac{2}{\pi}+o(1))\log n$. It is ambiguous in \cite{Er61} whether Erd\H{o}s intended the coefficients to be uniformly chosen from $\{-1,1\}$ or $\{0,1\}$. In the latter case, the constant $\frac{2}{\pi}$ should be $\frac{1}{\pi}$ (see the discussion in the comments). In the case of $\{-1,1\}$ Do \cite{Do24} proved that, if $R_n[-1,1]$ counts the number of roots in $[-1,1]$, then, almost surely,\[\lim_{n\to \infty}\frac{R_n[-1,1]}{\log n}=\frac{1}{\pi}.\]See also [522] .




References

[Do24] Y. Do, A strong law of large numbers for real roots of random polynomials . arXiv:2403.06353 (2024).

[EO56] Erd\H{o}s, Paul and Offord, A. C., On the number of real roots of a random algebraic equation . Proc. London Math. Soc. (3) (1956), 139-160.

[Er61] Erd\H{o}s, Paul, Some unsolved problems . Magyar Tud. Akad. Mat. Kutat\'{o} Int. K\"{o}zl. (1961), 221-254.


\bigskip
\noindent\small{Source: \url{https://www.erdosproblems.com/521}}
\end{document}
